\documentclass[a4paper,11pt]{article}
\usepackage[utf8]{inputenc}
\usepackage[T1]{fontenc}
\usepackage[french]{babel}
\usepackage{lmodern}
\usepackage{amsmath,amssymb,amsthm}
\usepackage{mathrsfs}
\usepackage{enumitem}
\usepackage{multicol}

\setlist[enumerate]{label=\textbf{\textcolor{Couleur8}{\arabic*.}}, left=1em}

% Si vous voulez aussi modifier les listes enumerate imbriquées :
\setlist[enumerate,2]{label=\textcolor{Couleur8}{\alph*)}}
\setlist[enumerate,3]{label=\textcolor{Couleur8}{\roman*)}}
\setlist[enumerate,4]{label=\textcolor{Couleur8}{\Alph*)}}
\usepackage{graphicx}
\usepackage{xcolor}
\usepackage{tcolorbox}
\usepackage{geometry}
\usepackage{fancyhdr}
\usepackage{titlesec}
\usepackage{cancel}
\usepackage{ifthen}
\usepackage{tikz}
\usetikzlibrary{arrows.meta}
\usepackage{tkz-tab}
\usepackage{eurosym}

% OPTION PROF/ÉLÈVE
% Décommentez UNE des deux lignes suivantes :
\newboolean{prof}\setboolean{prof}{true}   % Version professeur (avec corrections des devoirs)
%\newboolean{prof}\setboolean{prof}{false}  % Version élève (sans corrections des devoirs)

\definecolor{Couleur1}{HTML}{4A5D7A} %Th
\definecolor{Couleur2}{HTML}{A5B5D6} %Prop
\definecolor{Couleur3}{HTML}{A8C68A} %Méthode
\definecolor{Couleur6}{HTML}{8A9CC1} %Def
\definecolor{Couleur7}{HTML}{E6B459} %Rq Voc
\definecolor{Couleur8}{HTML}{D36740} %Titres
\definecolor{correction}{HTML}{D36740} % Couleur pour les corrections
\definecolor{coopmathscorr}{HTML}{F15929} % Couleur corrections coopmaths

% Configuration des marges
\geometry{
	top=2cm,
	bottom=2cm,
	inner=1.5cm,
	outer=1.5cm,
}

% Police
\renewcommand{\familydefault}{\sfdefault}

% Compteurs
\newcounter{seancenum}
\newcounter{seancenumD}
\setcounter{seancenum}{1}
\setcounter{seancenumD}{1}

% Configuration des en-têtes et pieds de page
\pagestyle{fancy}
\fancyhf{}
\ifthenelse{\boolean{prof}}{
    \fancyhead[L]{\textcolor{Couleur8}{\textbf{Corrections Automatismes - Classe de seconde - Version Professeur}}}
}{
    \fancyhead[L]{\textcolor{Couleur8}{\textbf{Corrections Automatismes - Séance 24 - Classe de seconde}}}
}
\fancyhead[R]{\textcolor{Couleur8}{\textbf{2025-2026}}}
\fancyfoot[L]{\textcolor{Couleur8}{Mme Bertail et Mme Pinto}}
\fancyfoot[R]{\textcolor{Couleur8}{\thepage}}
\renewcommand{\headrulewidth}{1pt}
\renewcommand{\headrule}{\hbox to\headwidth{\color{Couleur8}\leaders\hrule height \headrulewidth\hfill}}

% Environnements personnalisés
\newtcolorbox{seance}[1]{
    colback=Couleur6!10,
    colframe=Couleur1!65,
    fonttitle=\bfseries\large,
    title=Correction Séance \theseancenum #1,
    left=5mm,
    right=5mm,
    top=3mm,
    bottom=3mm,
    arc=0mm,
    after=\stepcounter{seancenum}
}

\newtcolorbox{seanceD}[1]{
    colback=Couleur6!5,
    colframe=Couleur6!45,
    fonttitle=\bfseries\large,
    title=Correction Devoirs - Séance \theseancenumD #1,
    left=5mm,
    right=5mm,
    top=3mm,
    bottom=3mm,
    arc=0mm,
    after=\stepcounter{seancenumD}
}

\begin{document}

\setcounter{seancenum}{24}
\newpage
\begin{seance}{} %24

\begin{enumerate}
\item Les solutions sont les abscisses des points d'intersection entre les deux courbes :
   $S=\{{\color[HTML]{f15929}\boldsymbol{-1\,;\,2}}\}$. 
	\item En mettant les fractions au même dénominateur, on a bien $\dfrac{3}{5}+\dfrac{1}{10}=\dfrac{7}{10}$.
    Réponse : {\bfseries \color[HTML]{f15929}Vrai}
	\item On sait d'après le cours que le coefficient directeur $m$ est donné par : $m=\dfrac{y_L-y_K}{x_-x_A}$.\\On applique avec les données de l'énoncé :
        $m=\dfrac{2-1}{-1-10}=
        {\color[HTML]{f15929}\boldsymbol{\dfrac{-1}{11}}}$.
	\item Comme $1{,}28=1+0{,}28=1+\dfrac{28}{100}$, multiplier par $1{,}28$ revient à augmenter de ${\color[HTML]{f15929}\boldsymbol{28}}\,\%$. 
	\item Augmenter de $100\,\%$ revient à multiplier par $2$. Augmenter de $60\,\%$ revient à multiplier par $1{,}6$.\\
        Le coefficient multiplicateur global est donc $2\times 1{,}6=3{,}2$.\\
        Ainsi, le taux global est donné par $3{,}2-1={\color[HTML]{f15929}\boldsymbol{220}}\,\%$.
	\item $f(-5)=\left(-5 \right)^2 +5 \times \left(-5\right) +1 = 1$. 
    On a donc $f(-5)={\color[HTML]{f15929}\boldsymbol{1}}$.
	\item On calcule la moyenne $m$ en faisant le quotient de la somme des notes coefficientées par la somme des coefficients. \\
	\medskip
	
        $m=\dfrac{12+10+2\times 9}{1+1+2}=\dfrac{40}{4}={\color[HTML]{f15929}\boldsymbol{10}}$
	\item La relation entre la probabilité d'un événement $A$ et celle de son contraire $\overline{A}$ est :  $P(\overline{A})=1-P(A)$.\\
          Ainsi : $P(\overline{A})=1-\dfrac{3}{10}={\color[HTML]{f15929}\boldsymbol{\dfrac{7}{10}}}$.
	\item Le vecteur d'origine $B$ égal au  vecteur $\overrightarrow{AB}$ est le vecteur $\overrightarrow{BG}$. 
    Ainsi, $\overrightarrow{AB}=\overrightarrow{B{\color[HTML]{f15929}\boldsymbol{G}}}$.
	
	\item $10$ stylos coûtent $15$\,\euro{} donc $5$ stylos coûtent  $7{,}50$\,\euro{}.Ainsi,   $15$ stylos coûtent $15+ 7{,}50 ={\color[HTML]{f15929}\boldsymbol{22{,}50}}$\,\euro{}.
	\item Le coefficient directeur $m$ de la droite $(AB)$ est donné par : $m=\dfrac{y_B-y_A}{x_B-x_A}=\dfrac{-1-(-3)}{3-0}=\dfrac{{\color{blue}\boldsymbol{2}}}{{\color{red}\boldsymbol{3}}}{\color[HTML]{f15929}\boldsymbol{}}$.

\medskip
\begin{tikzpicture}[baseline,scale = 0.6]

    \tikzset{
      point/.style={
        thick,
        draw,
        cross out,
        inner sep=0pt,
        minimum width=5pt,
        minimum height=5pt,
      },
    }
    \clip (-5,-5.25) rectangle (5,1);
    	\draw[color={blue},line width = 2] (-41.60251471689219,-30.735009811261456)--(44.60251471689219,26.735009811261456);
	\draw[color ={black},line width = 1.2] (-5,0)--(5,0);
	\draw[color ={black},line width = 1.2] (0,-5)--(0,5);
	\draw[color ={black},opacity = 0.4] (-5,1)--(5,1);
	\draw[color ={black},opacity = 0.4] (-5,-1)--(5,-1);
	\draw[color ={black},opacity = 0.4] (-5,2)--(5,2);
	\draw[color ={black},opacity = 0.4] (-5,-2)--(5,-2);
	\draw[color ={black},opacity = 0.4] (-5,3)--(5,3);
	\draw[color ={black},opacity = 0.4] (-5,-3)--(5,-3);
	\draw[color ={black},opacity = 0.4] (-5,4)--(5,4);
	\draw[color ={black},opacity = 0.4] (-5,-4)--(5,-4);
	\draw[color ={black},opacity = 0.4] (-5,5)--(5,5);
	\draw[color ={black},opacity = 0.4] (-5,-5)--(5,-5);
	\draw[color ={black},opacity = 0.4] (1,-5)--(1,5);
	\draw[color ={black},opacity = 0.4] (-1,-5)--(-1,5);
	\draw[color ={black},opacity = 0.4] (2,-5)--(2,5);
	\draw[color ={black},opacity = 0.4] (-2,-5)--(-2,5);
	\draw[color ={black},opacity = 0.4] (3,-5)--(3,5);
	\draw[color ={black},opacity = 0.4] (-3,-5)--(-3,5);
	\draw[color ={black},opacity = 0.4] (4,-5)--(4,5);
	\draw[color ={black},opacity = 0.4] (-4,-5)--(-4,5);
	\draw[color ={black},opacity = 0.4] (5,-5)--(5,5);
	\draw[color ={black},opacity = 0.4] (-5,-5)--(-5,5);
	\draw[color ={black},line width = 1.2] (0,-0.1)--(0,0.1);
	\draw[color ={black},line width = 1.2] (1,-0.1)--(1,0.1);
	\draw[color ={black},line width = 1.2] (-1,-0.1)--(-1,0.1);
	\draw[color ={black},line width = 1.2] (2,-0.1)--(2,0.1);
	\draw[color ={black},line width = 1.2] (-2,-0.1)--(-2,0.1);
	\draw[color ={black},line width = 1.2] (3,-0.1)--(3,0.1);
	\draw[color ={black},line width = 1.2] (-3,-0.1)--(-3,0.1);
	\draw[color ={black},line width = 1.2] (-0.1,0)--(0.1,0);
	\draw[color ={black},line width = 1.2] (-0.1,1)--(0.1,1);
	\draw[color ={black},line width = 1.2] (-0.1,-1)--(0.1,-1);
	\draw[color ={black},line width = 1.2] (-0.1,2)--(0.1,2);
	\draw[color ={black},line width = 1.2] (-0.1,-2)--(0.1,-2);
	\draw[color ={black},line width = 1.2] (-0.1,3)--(0.1,3);
	\draw[color ={black},line width = 1.2] (-0.1,-3)--(0.1,-3);
	\draw[opacity=0.8] (1,-0.4) node[anchor = center, rotate=0] {\scriptsize \color{black}$1$};
	\draw[opacity=0.8] (2,-0.4) node[anchor = center, rotate=0] {\scriptsize \color{black}$2$};
	\draw[opacity=0.8] (3,-0.4) node[anchor = center, rotate=0] {\scriptsize \color{black}$3$};
	\draw[opacity=0.8] (4,-0.4) node[anchor = center, rotate=0] {\scriptsize \color{black}$4$};
	\draw[opacity=0.8] (-1,-0.4) node[anchor = center, rotate=0] {\scriptsize \color{black}$-1$};
	\draw[opacity=0.8] (-2,-0.4) node[anchor = center, rotate=0] {\scriptsize \color{black}$-2$};
	\draw[opacity=0.8] (-3,-0.4) node[anchor = center, rotate=0] {\scriptsize \color{black}$-3$};
	\draw[opacity=0.8] (-4,-0.4) node[anchor = center, rotate=0] {\scriptsize \color{black}$-4$};
	\draw[opacity=0.8] (-0.6,1.1) node[anchor = center, rotate=0] {\scriptsize \color{black}$1$};
	\draw[opacity=0.8] (-0.6,2.1) node[anchor = center, rotate=0] {\scriptsize \color{black}$2$};
	\draw[opacity=0.8] (-0.6,3.1) node[anchor = center, rotate=0] {\scriptsize \color{black}$3$};
	\draw[opacity=0.8] (-0.6,4.1) node[anchor = center, rotate=0] {\scriptsize \color{black}$4$};
	\draw[opacity=0.8] (-0.6,-0.9) node[anchor = center, rotate=0] {\scriptsize \color{black}$-1$};
	\draw[opacity=0.8] (-0.6,-1.9) node[anchor = center, rotate=0] {\scriptsize \color{black}$-2$};
	\draw[opacity=0.8] (-0.6,-2.9) node[anchor = center, rotate=0] {\scriptsize \color{black}$-3$};
	\draw[opacity=0.8] (-0.6,-3.9) node[anchor = center, rotate=0] {\scriptsize \color{black}$-4$};
	\draw[color ={black},line width = 1.25,opacity = 0.8] (-0.13,-2.88)--(0.13,-3.13);\draw[color ={black},line width = 1.25,opacity = 0.8] (-0.13,-3.13)--(0.13,-2.88);
	\draw[opacity=1] (0.1,-3.2) node[anchor = center, rotate=0] {\normalsize \color{black}$A$};
	\draw[opacity=1] (3,-0.5) node[anchor = center, rotate=0] {\normalsize \color{black}$B$};
	\draw[color ={black},line width = 1.25,opacity = 0.8] (2.88,-0.88)--(3.13,-1.13);\draw[color ={black},line width = 1.25,opacity = 0.8] (2.88,-1.13)--(3.13,-0.88);
	\draw[opacity=1] (-0.2,-0.3) node[anchor = center, rotate=0] {\scriptsize \color{black}$\text{O}$};
	\draw[color ={black},line width = 2, dashed ] (0,-3)--(3,-3);
	\draw[color ={black},line width = 2, dashed ] (3,-1)--(3,-3);
	\draw[opacity=1] (1.5,-3.3) node[anchor = center, rotate=0] {\normalsize \color{red}$3$};
	\draw[opacity=1] (3.5,-2) node[anchor = center, rotate=0] {\normalsize \color{blue}$2$};

\end{tikzpicture}
	\item L'antécédent de $8$ est le nombre $x$ qui a pour image $8$. On cherche donc $x$ tel que : $2x-2=8$ \\Soit $x=\dfrac{8+2}{2}={\color[HTML]{f15929}\boldsymbol{5}}$.
	\item  On reconnaît une équation du type $x^2=k$ avec $k=-11$.\\Puisque $-11$ est strictement négatif, l'équation n'a pas de solution. Ainsi, $S={\color[HTML]{f15929}\boldsymbol{\emptyset}}$.
	\item On utilise l'égalité remarquable $(a-b)^2=a^2-2ab+b^2$ avec $a=5$ et $b=3x$.\\$
       (5-3x)^2=5^2-2\times 5\times 3x+ (3x)^2={\color[HTML]{f15929}\boldsymbol{9x^2-30x+25}}$
	\item On utilise le triangle rectangle représenté ci-dessous et on applique le théorème de Pythagore : \\
    $AB^2=AC^2+BC^2$ donc 
         $AB^2= 4^2+3^2$ donc 
         $AB=\sqrt{25}$ donc 
$AB=5$

             La valeur exacte de $AB$ est ${\color[HTML]{f15929}\boldsymbol{\sqrt{25}}}$ u.l.\begin{tikzpicture}[baseline,scale = 0.5]

    \tikzset{
      point/.style={
        thick,
        draw,
        cross out,
        inner sep=0pt,
        minimum width=5pt,
        minimum height=5pt,
      },
    }
    \clip (-1,0) rectangle (8,5.3);
    	\draw[color ={gray}] (0,0)--(0,5);
	\draw[color ={gray}] (1,0)--(1,5);
	\draw[color ={gray}] (2,0)--(2,5);
	\draw[color ={gray}] (3,0)--(3,5);
	\draw[color ={gray}] (4,0)--(4,5);
	\draw[color ={gray}] (5,0)--(5,5);
	\draw[color ={gray}] (6,0)--(6,5);
	\draw[color ={gray}] (7,0)--(7,5);
	\draw[color ={gray}] (8,0)--(8,5);
	\draw[color ={gray}] (0,0)--(8,0);
	\draw[color ={gray}] (0,1)--(8,1);
	\draw[color ={gray}] (0,2)--(8,2);
	\draw[color ={gray}] (0,3)--(8,3);
	\draw[color ={gray}] (0,4)--(8,4);
	\draw[color ={gray}] (0,5)--(8,5);
	\draw[color ={blue},line width = 3] (1,4)--(5,1);
	\draw[color ={black},line width = 2] (6,4)--(7,4);
	\draw[color ={blue},line width = 2, dashed ] (1,4)--(5,4);
	\draw[color ={blue},line width = 2, dashed ] (5,4)--(5,1);
	\draw[opacity=1] (6.5,4.5) node[anchor = center, rotate=0] {\scriptsize \color{black}$1 \text{u.l.}$};
	\draw[color ={black},line width = 0.625,opacity = 0.8] (0.81,4.19)--(1.19,3.81);\draw[color ={black},line width = 0.625,opacity = 0.8] (0.81,3.81)--(1.19,4.19);\draw[color ={black},line width = 0.625,opacity = 0.8] (4.81,1.19)--(5.19,0.81);\draw[color ={black},line width = 0.625,opacity = 0.8] (4.81,0.81)--(5.19,1.19);\draw[color ={black},line width = 0.625,opacity = 0.8] (4.81,4.19)--(5.19,3.81);\draw[color ={black},line width = 0.625,opacity = 0.8] (4.81,3.81)--(5.19,4.19);
	\draw [color={black}] (1,4.5) node[anchor = center,scale=1, rotate = 0] {A};
	\draw [color={black}] (5,0.5) node[anchor = center,scale=1, rotate = 0] {B};
	\draw [color={black}] (5,4.5) node[anchor = center,scale=1, rotate = 0] {C};

\end{tikzpicture}
\end{enumerate}

\end{seance}

\end{document}